% =============================================================================
% CHAPTER 1: Introduction
% =============================================================================
% SPDX-FileCopyrightText: 2026 Frank Langbein <frank@langbein.org>
% SPDX-License-Identifier: LPPL-1.3c
%
% Suggested sections: Motivation and context; Aims and objectives; Contributions;
% Dissertation structure (roadmap). Optional: Scope (boundaries of the work),
% Terminology or key terms. Start with a short chapter intro (no \section) before
% the first section. This file also demonstrates template features; see
% Section~\ref{sec:ch1-reference}. Remove that section for your own report.
% =============================================================================

% -----------------------------------------------------------------------------
% Chapter intro (no section): one short paragraph before the first \section
% -----------------------------------------------------------------------------
\newthought{Media Players} are devices that focus on the delivery of music files, predominantly for leisure. Having a fondness for music, the Author aimed to make the playback of music as accessible as possible for the cognitively impaired. This project has a particular focus on patients with dementia, an umbrella term of diseases that impact cognitive ability.

% -----------------------------------------------------------------------------
% Motivation and context (typical first section)
% -----------------------------------------------------------------------------
% Use \section, \subsection, \label{key}. Refer with \ref{key} or \autoref{key}.
\section{Motivation and context}\label{sec:motivation}
Dementia is a syndrome that the Author has had experience with before through meeting an extended family member. Seeing the gravity of the mental degradation that can occur, when the Author saw the opportunity to develop a solution for the betterment of the 'quality of life' for these patients, they immediately expressed interest. Being the first project the Author has done in the field of accessibility, the Author had much to learn and they were excited to do so.

Doing preliminary research into media players, there did exist tailored devices for the patients with dementia on the market already but none of these contained any form of remote management, which was a gap in the market the Author decided to address.

% Aims and objectives (typical second section)
% -----------------------------------------------------------------------------
% Use \citep{} (parenthetical) or \citet{} (narrative), not \cite.

\section{Aims and objectives}\label{sec:objectives}
The aims and objectives of the project were as follows:
\begin{enumerate}
  \item Develop a tangible device that can be interacted with remotely which is able to play given media.
  \item Engage with social care researchers in a productive manner throughout the development of the device.
  \item Ensure the device is simple to use, lowering the barrier of entry for the device.
  \item Develop a secondary way of interaction with the device if time allows it.
\end{enumerate}

% -----------------------------------------------------------------------------
% Dissertation structure / roadmap (typical final section of introduction)
% -----------------------------------------------------------------------------

\section{Dissertation structure}\label{sec:roadmap}
The Dissertation you are about to read will have the following structure: 
\begin{itemize}
  \item \textbf{Chapter~\ref{ch2}} (Background): context, related work, problem statement.
  \item \textbf{Chapter~\ref{ch3}} (Methodology): research design, specification, system design, implementation.
  \item \textbf{Chapter~\ref{ch4}} (Results and Evaluation): setup, results, evaluation, limitations.
  \item \textbf{Chapter~\ref{ch5}} (Conclusions and Future Work): summary and next steps.
  \item \textbf{Chapter~\ref{ch6}} (Reflection on Learning): transferable learning.
  \item \textbf{Appendix~\ref{app:results}}: additional results.
\end{itemize}

